\subsection{Die Wand-Lesart: gemeinsamer geometrischer Anschluss und korrigierte Neutrino-Weiche}\label{sec:wallaudit}
\key{\textbf{Die Wand-Idee liefert einen bedingten geometrischen Anschluss, aber keinen Zwang zu Dirac-Neutrinos.} Die volle ursprüngliche Quelle enthält bereits Standardmodell-Singuletts mit $B-L=\pm2$. Ein spezielles Spinorbündel verbindet außerdem Familienindex und Anomalieklasse. Die entscheidende gemeinsame Auswahl betrifft die Verbindung, den physischen Zustand und die Kopplung; sie folgt nicht aus dem Fehlen eines Feldes in der ersten 248-dimensionalen Stufe.}

Das am 29.~September erneut gelesene Claude-Artefakt und der beigefügte Wand-Bericht werden hier gegen die Primärliteratur und den nativen TFPT-Code geprüft. Die dort genannten 27 externen Prüfungen und die nicht vorliegenden Commitinhalte werden nicht als eigene Ausführung verbucht. Die Aussage über die vollständige Quelle lässt sich dagegen unmittelbar im vorhandenen Quellprogramm entscheiden.

\paragraph{Was Hořava--Witten tatsächlich verbindet.}
Die Konstruktion beginnt mit elfdimensionaler Supergravitation und zwei zehndimensionalen Randkomponenten. Deren chirale Anomalie verlangt unter diesen Voraussetzungen 248 Vektormultipletts pro Wand; die vollständige Faktorisierung und das heterotische Wörterbuch führen zu $E_8$~\cite{horavawittenorigin,horavawittenwall}. Der relevante irreduzible Gravitationsanteil ist $\Tr R^6$ in einem Zwölfform-Anomaliepolynom. Eine Quartikidentität allein ersetzt diese Rechnung nicht. Die zweidimensionale Quellenanomalie $c_-=8$ und dieses zehn-dimensionale Inflow sind verschiedene Aussagen, solange ihre geometrische Abbildung fehlt.

Die bereits bekannte Verzweigung
\begin{equation}
248=(45,1)\oplus(1,15)\oplus(16,4)
\oplus(\overline{16},\overline4)\oplus(10,6)
\end{equation}
passt zur heterotischen Wahl eines $SU(4)$-Eichbündels mit Kommutant $Spin(10)$. Das Eich-$SU(4)$ wirkt auf Bündelfasern. Das isomorphe Tangential-$Spin(6)$ wirkt auf innere Raumrichtungen und deren Spinoren. Eine Gleichsetzung dieser Wirkungen ist eine zusätzliche, prüfbare Bündelwahl. Auch
\begin{equation}
d=10-2(8-g)=2g-6
\end{equation}
ist nur innerhalb der hinzugefügten zehn-dimensionalen Wand- und Tangentialeinbettung gültig. Sie leitet weder diese Wand noch die Einbettung aus P1/P2 her. Die vierdimensionale $\Z_{16}$-Anomalie aus Abschnitt~\ref{sec:spinliftsynthesis} darf dabei nicht unverändert für Kandidaten anderer Raumzeitdimension verwendet werden: Ihr Eingang ist bereits ein vierdimensionales fermionisches Symmetrieproblem.

Auch die Familienplätze müssen wirklich übersetzt werden. In einer heterotischen Entwicklung
\begin{equation}
\lambda^{aA}(x,y)=\sum_i\psi_i^a(x)\,\xi_i^A(y)
\end{equation}
bezeichnet $A=1,\ldots,4$ den $SU(4)$-Bündelindex, während $i$ die unabhängigen inneren Nullmoden und damit die vierdimensionalen Familien zählt. Die drei beibehaltenen TFPT-Plätze $C_1,C_2,C_3$ sind deshalb nicht bereits durch das gemeinsame Zeichen $(16,4)$ mit drei solchen Nullmoden identifiziert. Dafür braucht es einen Intertwiner, der insbesondere die ursprüngliche Familienmonodromie erhält.

Ebenso wenig folgt die Spin-Ladungsrelation automatisch für das volle masselose heterotische Spektrum. Ein ungebrochenes $\mathcal N=1$-chirales Multiplett enthält Skalar und Fermion in derselben Eichdarstellung. Ist $\omega$ nur das entsprechende Eichzentrum, wirkt $\omega^2$ auf beide gleich, $(-1)^F$ dagegen verschieden. Eine zusätzliche diagonale Symmetrie, Projektion oder SUSY-Brechung muss ausdrücklich definiert werden, bevor die Relation als physische Spektralbedingung verwendet wird.

\paragraph{Die stärkste richtige geometrische Identität.}
Sei $X$ eine geschlossene orientierte Spin-Sechsmannigfaltigkeit, und werde ausdrücklich $V=S_X^+$ gewählt. Für die formalen Krümmungswurzeln $x_1,x_2,x_3$ von $TX$ sind die vier Spinorgewichte
\begin{equation}
\frac12(x_1+x_2+x_3),\quad\frac12(x_1-x_2-x_3),\quad
\frac12(-x_1+x_2-x_3),\quad\frac12(-x_1-x_2+x_3).
\end{equation}
Die elementarsymmetrischen Polynome ergeben unmittelbar
\begin{equation}\label{eq:wallchern}
c_1(S_X^+)=0,\qquad c_2(S_X^+)=-\frac12p_1(TX),\qquad
c_3(S_X^+)=e(TX),
\end{equation}
mit umgekehrtem Vorzeichen des letzten Terms bei vertauschter Chiralität. Da $\widehat A$ keinen Term vom Grad sechs und $c_1(V)=0$ hat, folgt
\begin{equation}\label{eq:wallindex}
\operatorname{ind}D_{S_X^+}
=\int_X\widehat A(TX)\operatorname{ch}(S_X^+)
=\frac12\int_X c_3(S_X^+)=\frac{\chi(X)}2.
\end{equation}
Diese spezielle Eulerformel ist also korrekt. Für ein allgemeines heterotisches $SU(4)$-Bündel gilt dagegen $\operatorname{ind}D_V=\tfrac12\int_Xc_3(V)$; es ist keine Eulerformel für $TX$. Ein Index zählt Nettochiralität, nicht automatisch die Abwesenheit zusätzlicher vektorartiger Paare. Ein konkreter Gegenvergleich ist das heterotische Modell von Braun et al.: Dort hat $X$ die Hodgezahlen $(3,3)$, also $\chi(X)=0$, während $c_3(\widetilde V)=\pm54$ auf der neunfachen Überlagerung drei Nettofamilien auf $X$ liefert~\cite{braunheterotic}.

Auf einer Calabi--Yau-Dreifaltigkeit wird auch $c_2(S_X^+)=c_2(TX)$: Aus $c_1(TX)=0$ folgt $p_1(TX)=-2c_2(TX)$. Damit passen Familienindex und die topologische Bianchi-Bedingung
\begin{equation}
c_2(TX)-c_2(V)-c_2(V_{\rm hidden})=[W]
\end{equation}
mit trivialem verborgenem Bündel und $[W]=0$ zusammen. Dies löst die \emph{kohomologische} Bedingung; eine lokale Lösung der gesamten Feldgleichungen folgt daraus noch nicht.

Die Verbindung bleibt entscheidend: Bei Calabi--Yau-Holonomie $SU(3)$ zerfällt $S_X^+$ als $\mathbf1\oplus\mathbf3$, bis auf Konjugation. Die Levi-Civita-Verbindung in diesem Bündel lässt $E_6$ statt $Spin(10)$ übrig. Das wurde zusätzlich an den ursprünglichen 240 E8-Wurzeln nachgerechnet: Orthogonalität zum tatsächlichen $A_2$ liefert 72 Wurzeln plus sechs Cartanrichtungen, also Dimension 78; für das volle $A_3$ sind es 40 plus fünf, also Dimension 45. Eine gleiche abstrakte Strukturgruppe legt ihre Verbindung und Holonomie nicht fest.

\paragraph{Ein konkreter Anschluss, ohne eine neue Innengeometrie zu erfinden.}
Braun--Candelas--Davies--Donagi verwenden einen bekannten Raum mit
$(h^{1,1},h^{2,1})=(1,4)$, $\chi=-6$ und $\pi_1=\Z_{12}$. Ihr Abschnitt~3 untersucht Deformationen von $TX\oplus\mathcal O^{\oplus r}$; Anhang~B skizziert die mögliche Rang-vier-Konstruktion als Hilfsargument; das ausgearbeitete Modell verwendet Rang fünf. Die Topologie bleibt dabei erhalten~\cite{braundeformation}. Für die Wand-Lesart ist deshalb der präzise Kandidat eine stabile irreduzible, mit dem Quotienten verträgliche $SU(4)$-Deformation von $TX\oplus\mathcal O$. Sie könnte $c_2$ und den Drei-Familien-Index erhalten und zugleich den Kommutanten verkleinern. Ihre tatsächliche Verbindung, Nullmoden, Wilsonlinien und Auswahl aus TFPT müssen gemeinsam bestimmt werden. Volle Tangentialholonomie $Spin(6)$ darf dabei nicht zugleich als Calabi--Yau-Eigenschaft gefordert werden.

Die Deformationsrichtungen liegen unter den üblichen Calabi--Yau-Voraussetzungen in
\begin{equation}
H^1(\operatorname{End}(TX\oplus\mathcal O))=
H^1(\operatorname{End}TX)\oplus H^1(TX)\oplus H^1(T^*X).
\end{equation}
Die beiden letzten Summanden mischen den Dreier- und den Singulettanteil. Integrabilität, Stabilität und Äquivarianz sind zusätzliche Bedingungen, keine Folgen bloß nichtverschwindender Kohomologie. Eine solche Bündeldeformation wird physikalisch als Higgs-Hintergrund gelesen; sie ist noch nicht mit dem nativen symmetrischen $B_{ab}$-Kanal identifiziert. Genau hier wäre eine konkrete Gleichsetzung von Wirkungen statt einer Dimensionszählung erforderlich.

Dafür existiert ein konkreter mathematischer Fortsetzungssatz: Li--Yau zeigen für eine gegebene kleine holomorphe Deformationsfamilie einer Summe zweier stabiler Grad-null-Bündel, dass bei beidseitig nichtverschwindenden gemischten Kodaira--Spencer-Komponenten eine passende Hermitian--Yang--Mills-Metrik fortgesetzt werden kann. Ihr Satz~4.3 behandelt unter weiteren Voraussetzungen das gekoppelte Strominger-System~\cite{liyautorsion}. Das ist ein vorhandenes Konstruktionswerkzeug für die bedingte Wand-Realisierung; die Deformationsfamilie selbst und ihr TFPT-Ursprung sind weiterhin Eingang.

\paragraph{Die Herleitung eines Ausschlusses aus der 248 ist widerlegt.}
Wir verwenden die bereits originale Normierung
\begin{equation}
Q_{\rm em}=T_3+Y,\qquad X=5(B-L)-4Y,\qquad \omega=i^X.
\end{equation}
Für die linkshändigen $\nu^c$ gilt $(X,Y,B-L)=(5,0,1)$. Die zwei exemplarischen nativen Zustände
\begin{align}
B_{11}&=C_1(-3)C_1(-1)\Omega,& \norm{B_{11}}^2&=1,\\
B_{12}&=\bigl(C_1(-3)C_2(-1)+C_2(-3)C_1(-1)\bigr)\Omega,
&\norm{B_{12}}^2&=2
\end{align}
haben $X=10$, $Y=0$, $B-L=2$. Alle acht positiven und negativen Wurzelgeneratoren von $SU(3)_c\times SU(2)_L$ annihilieren sie im einfachen unitären Quotienten; auch ihre Cartan-Gewichte verschwinden. Dieselbe Rechnung wurde für alle zehn symmetrischen Paare und alle zehn adjungierten Gegenladungen ausgeführt. Sie sind damit wirkliche SM-Singuletts der Quelle, keine bloß passend benannten Vektoren.

Die Spin-Ladungsrelation lässt sie zu:
\begin{equation}
X=\pm10\equiv2\pmod4,\qquad \omega=-1,\qquad
\omega^2=+1=(-1)^F\quad\text{für einen Skalar}.
\end{equation}
Sie verbietet einen \emph{ungerade} geladenen Skalar wie den zuvor geprüften $\overline{16}_H$ unter dem zusätzlichen Spin-$\Z_4$-Postulat. Sie verbietet keinen bosonischen Ladung-zwei-Paarkanal. Die chiral bosonische Quelle allein beweist weiterhin keine vierdimensionale Skalarrealisierung.

Auch die niedrigste mögliche Stufe lässt sich ohne endliche Abschneideannahme bestimmen. Ein SM-Singulettmomentum mit $Y=0$ hat in den ersten fünf E8-Koordinaten denselben Wert $t$. Aus $X=-10t=10$ folgt $t=-1$. Die übrigen drei Koordinaten sind dann ganzzahlig und haben wegen der geraden E8-Gesamtsumme eine ungerade Summe. Also gilt
\begin{equation}
\norm{p}^2\ge5+1=6,\qquad h=\frac{\norm p^2}{2}+N\ge3.
\end{equation}
Die antisymmetrischen Gewicht-drei-Zustände erreichen diese Schranke; die symmetrische Familien-$10$ tritt auf Gewicht vier auf, wie in Abschnitt~\ref{sec:sourcepairs} bewiesen. Das Fehlen in Gewicht eins ist richtig, seine Übertragung auf die vollständige Quelle ist falsch.

\paragraph{Was über Dirac oder Majorana wirklich entscheidet.}
Bei einer \emph{exakt erhaltenen} schützenden Rest-$\Z_4$ lauten die Auswahlregeln für hyperladungsneutrale Operatoren:
\begin{center}
\begin{tabular}{lrl}
\toprule Operator & $X$ & Rest-$\Z_4$\\\midrule
$L H_u\nu^c$ & $0$ & erlaubt\\
$\nu^c\nu^c$ & $10$ & verboten\\
$(L H_u)^2$ & $-10$ & verboten\\
$(\nu^c)^4$ & $20$ & erlaubt\\
$(u^cd^cd^c)^2$ & $-10$ & verboten\\
$B^\dagger_{ab}\nu^c_a\nu^c_b$ & $0$ & erlaubt\\\bottomrule
\end{tabular}
\end{center}
Hier bedeutet erlaubt ausschließlich: von dieser Symmetrie nicht ausgeschlossen. Das bestimmt weder eine vierdimensionale Kopplung noch eine Zerfallsrate. Die elektroschwache Higgsphase erfordert die Eichkompensation
\begin{equation}
\widehat\omega=\omega e^{2\pi iY}
=e^{5\pi i(B-L)/2},\qquad \widehat\omega H_u=H_u.
\end{equation}
Auf hyperladungsneutralen Operatoren bleibt die gezeigte Auswahlregel gleich. Die globale Gruppenordnung auf quarkischen Bruchladungen hängt zusätzlich von den Eichzentrumsquotienten ab; aus dieser Formel allein wird keine neue globale Gruppe behauptet.

Heeck--Rodejohann enthalten ausdrücklich einen Skalar $\chi$ mit $B-L=-2$ und einen Skalar $\phi$ mit Ladung vier. Ihr Dirac-Zweig verlangt $\langle\chi\rangle=0$ bei nichtverschwindendem $\langle\phi\rangle$~\cite{heeckquadruple}. Gerade dieses zitierte Beispiel widerlegt die Gleichsetzung von Feldabwesenheit und Dirac-Schutz. Ein nachgewiesener Neutron--Antineutron-Übergang würde den hier beschriebenen exakten Restgruppen-Zweig ausschließen, nicht jede denkbare Theorie mit Dirac-Neutrinos.

\paragraph{Warum der vorhandene Familienprozess die Phase noch nicht auswählt.}
Alle zwölf ursprünglichen $A_3$-Wurzeln besitzen $X=0$; seine Cartanoperatoren sind ebenfalls neutral. Der bisherige Familientransport und das neutrale $R$ erhalten diese Ladung. Im invarianten Quellvakuum folgt deshalb exakt
\begin{equation}
\langle B_{ab}\rangle_\Omega
=\langle\omega B_{ab}\omega^{-1}\rangle_\Omega
=-\langle B_{ab}\rangle_\Omega=0,
\end{equation}
während $\langle B^\dagger B\rangle$ positiv ist. Die volle Transportgleichung aus Abschnitt~\ref{sec:masssourceconnection} erhält entsprechend die Null-Lösung. Sie kann einen eingesetzten Massentensor transportieren, aber seine Erzeugung nicht aus einer homogenen Gleichung ableiten. Die Wahl einer anderen Phase oder einer geladenen Randpräparation benötigt die zugehörige ursprüngliche Wirkung. In einer Eichtheorie ist ein geladener Einpunktwert zudem nur nach Eichfixierung ein direktes Ordnungsparameterbild; eine eichinvariante Higgsdiagnose verwendet passende Korrelationen. Aus dem Vakuum-Einpunktwert wird hier kein Verbot einer Higgsphase gefolgert.

Ein neutraler Skalaron ist vor Symmetriebrechung ebenfalls nicht dasselbe Feld wie ein komplexer Skalar mit $B-L=2$. Eine mögliche spätere Verbindung wäre dessen neutrale Radialmode in einer Darstellung $\sigma=(v+s)e^{ia/v}/\sqrt2$ samt hergeleiteter Mischung mit dem Skalaron. Der Ledger bezeichnet $\sigma=\text{Skalaron}$ als bedingt. Die entsprechende Zeile in \file{verification/v254_inner_fluctuations.py} prüft wörtlich den Wahrheitswert \texttt{True}; der frische Lauf mit fünf grünen Prüfungen beweist diese Identifikation daher nicht.

\paragraph{Die gemeinsame Bedingung für den weiteren Lösungsweg.}
Die Wand-Lesart darf die vorhandenen $\alpha$-, Massen- und Flavorherleitungen nicht durch freie Kompaktifizierungsparameter ersetzen. Schon dimensionsanalytisch ist $g_{10}^2$ von der Dimension Länge hoch sechs, während $c_3=1/(8\pi)$ dimensionslos ist. In einer unverwobenen Produktkompaktifizierung mit konstantem Hintergrund gilt schematisch $g_4^2=g_{10}^2/\operatorname{Vol}(X)$; Normierung, Schwellen und Warp-/Dilatonkorrekturen sind zusätzlich zu bestimmen. Der gleiche Gruppenname liefert diesen Zahlenanschluss nicht.

Der konsistente positive Auftrag ist deshalb enger und zusammenhängend: Aus dem ursprünglichen markierten Prozess muss ein gemeinsames physisches Wirkungs- und Zustandsfunktional hervorgehen. Seine Lösung muss den bereits vorhandenen Masseneinsatz im erlaubten geladenen Kanal tragen, den ursprünglichen Familientransport erhalten und zugleich den neutralen Nahtport reproduzieren. Eine heterotische Bündelrealisierung wäre daran zu messen: an der tatsächlichen Verbindung und ihrem Spektrum, an ihrem physischen Masseneinsatz und an den gleichen Kopplungen. Zuse-/Hypergraph-Regeln und der Universalraum bleiben mögliche Prozess- beziehungsweise Vergröberungsdarstellungen dieses Funktionals; sie wählen es nicht allein durch ihre Darstellung aus.

\picturetext{Die Wand ist eine mögliche Bauform für dasselbe Instrument. Im vollständigen Instrument ist das für Majorana-Massen benötigte Bauteil bereits vorhanden. Der Familienweg beschreibt, wie es mitbewegt wird. Gesucht ist jetzt die aus den Ursprungsregeln folgende Betriebsweise, die dieses Bauteil tatsächlich aktiviert und dabei gleichzeitig die schon berechneten Massen, Mischungen und die Feinstrukturkonstante liefert.}

\subsection{Vom Quellenfeld zum wirklichen Massenblock und zur gemeinsamen Vergröberung}\label{sec:masszeromode}
\key{\textbf{Der innere Operatoranschluss ist jetzt explizit geschlossen.}
Die zuvor konstruierten Paarfelder erzeugen durch ihre tatsächliche Modenwirkung
genau die symmetrischen Matrixeinträge des ursprünglichen schweren
Neutrinoblocks. Dieselbe Einsetzung liefert nach Eliminierung der schweren
Zustände die vorhandene Seesaw-Matrix, einen Determinantenbeitrag und eine
vollständige Gedächtnisantwort. Die Wahl der Einsetzung und ihr physischer
Ursprung bleiben dabei die benannten Voraussetzungen.}

\textbf{Die entscheidende Prüfung betrifft die Wirkung.}
Ein Zustand mit passenden Ladungen genügt noch nicht als Massenoperator.
Für ein homogenes Quellenfeld $v$ vom chiralen Gewicht $h$ bezeichnet
$o(v)=v_{(h-1)}$ seine gradbewahrende Mode~\cite{dongradical}.
Aus der ursprünglichen Stromalgebra folgt für
$u_{ab}=C_a(-3)C_b(-1)\Omega$
\begin{align}
Y(u_{ab},z)&=\frac12:\partial^2 C_a(z)C_b(z):,\\
o(u_{ab})&=\frac12\sum_{m\in\Z}(m+1)(m+2)
 :C_a(m)C_b(-m):.
\end{align}
Auf dem gesamten Stromraum $V_1$ bleiben genau die Terme $m=0,1$ übrig:
\begin{equation}\label{eq:massmodeexact}
\boxed{o(u_{ab})\big|_{V_1}=C_b(0)C_a(0)+3C_b(-1)C_a(1).}
\end{equation}
Der Term $m=-1$ besitzt Koeffizient null; die übrigen Terme verschwinden
durch Grad und Vakuumbedingung. Die Formel ist damit keine frei entworfene
Matrix, sondern eine Folge des vorhandenen Feldoperators.

Setze $|a\rangle=C_a(-1)\Omega$ und
$|\bar a\rangle=C_a^\dagger(-1)\Omega$. Mit den ursprünglichen
Gitterphasen ergibt die Rechnung für alle vier Familienplätze
\begin{align}
o(B_{aa})\big|_{V_1}&=|a\rangle\langle\bar a|,\\
o(B_{ab})\big|_{V_1}&=|a\rangle\langle\bar b|
                    +|b\rangle\langle\bar a|\quad(a<b),\\
o(B_{ab}^\dagger)\big|_{V_1}&=
  \bigl(o(B_{ab})\big|_{V_1}\bigr)^\dagger.
\end{align}
Diese Gleichheiten wurden als vollständige PBW-Ausgaben auf allen
248 Basisrichtungen geprüft, einschließlich der Cartanrichtungen.
Zusätzlich erklärt die Ladung das gesamte Stützgebiet: $X(B)=10$ kann im
$V_1$-Spektrum nur $X=-5$ mit $X=+5$ verbinden. Diese beiden Räume sind
genau die vier $C^\dagger$- beziehungsweise $C$-Zustände.
Für die sechs Paarfelder mit Indizes $1,2,3$ verbleibt deshalb ein
sechsdimensionaler Wirkungsraum; die übrigen 242 Richtungen werden annihiliert.
Auf höheren Quellengraden darf diese endliche Aussage nicht fortgeschrieben
werden. Die unteren Feldmoden erzeugen dort weitere Zustände.

Die Gewicht-zwei-Paarfelder haben auf $V_1$ die Wirkung null.
Die Gewicht-drei-Felder liefern stattdessen
$|a\rangle\langle\bar b|-|b\rangle\langle\bar a|$.
Damit ist die Symmetrie der Familienmatrix tatsächlich entscheidend:
Ein lokaler Majorana-Massenterm für gleichartige Weylfelder benötigt eine
symmetrische Familienform. Die Gewicht-vier-Felder liefern genau diese
Form. Die triviale Wirkung auf dem Vakuum-Top-Level $V_0$ widerspricht
dem nicht; die hier berechnete Wirkung liegt auf $V_1$, also jenseits der
gewöhnlichen Zhu-Algebra $A_0$~\cite{barronhigherzhu}.

\textbf{Der vollständige Block und die richtige Konjugation.}
Sei $K=K^{\mathsf T}$ eine schwere Massenform in der nativen Dreierbasis.
Der bereits benutzte Zustandskodierer lautet $\Gamma(\overline K)$.
In der Reihenfolge $(C^\dagger,C)$ folgt nun als Operatorgleichheit
\begin{equation}\label{eq:sourcemajoranaoperator}
D_K=o(\Gamma(\overline K))+o(\Gamma(\overline K))^\dagger
=\begin{pmatrix}0&K\\K^\dagger&0\end{pmatrix},\qquad
\Tr D_K^2=2\Tr K^\dagger K=2\norm{\Gamma(\overline K)}^2.
\end{equation}
Das Spektrum besteht aus den positiven und negativen Singularwerten von
$K$. Der Block ist selbstadjungiert, ungerade bezüglich der vorhandenen
Teilchen-/Antiteilchen-Chiralitätsgraduierung und verträglich mit der
reellen Struktur aus Tausch und komplexer Konjugation. Diese endlichen
Eigenschaften beweisen noch keine Raumzeitstatistik.

Die ursprüngliche Familienabbildung $S$ wird dabei beibehalten:
$K_C=S^{\mathsf T}K_{\rm original}S$ und
$W=\diag(\overline S,S)$ geben
\begin{equation}
W D_{K_C}W^\dagger=D_{K_{\rm original}}.
\end{equation}
Ein direkter Vergleich mit \file{verification/v252_full_finite_triple.py}
setzt diesen Quellenblock in den vorhandenen 96-dimensionalen Operator
ein. Er ersetzt exakt dessen Majoranablock; die gleichzeitig vorhandenen
komplexen Dirac-/Yukawablöcke bleiben gleich. Für den bisherigen bedingten
schweren Ansatz entsteht somit genau
$K_{\rm original}=M_R/M_3=\diag(\varepsilon/3,2\varepsilon/3,1)$.
Die Werte dieser Einsetzung werden hierbei aus dem Original übernommen.

Zwei unterschiedliche Wirkungen des neutralen Feldes müssen auseinanderbleiben:
\begin{equation}
R_0\Gamma(K)=4\Gamma(K),\qquad
R_0\big|_{\Span(C_{1,2,3},C_{1,2,3}^\dagger)}=I_6,
\qquad [R_0,D_K]\big|_{V_1}=0.
\end{equation}
Ein Gewicht-zwei-Feld wirkt nicht wie ein gewöhnlicher
Ladungsableitungsoperator. Das vertraute Vergleichsbeispiel ist
$L_0B=4B$, aber $[L_0,B_{(3)}]=0$.
Für gemischte Paare $B_{0i}$ gilt der letzte Kommutator gerade nicht,
weil der ausgelassene Platz $C_0$ den $R_0$-Eigenwert $-3$ trägt.
Auf dem behaltenen Block gilt außerdem
$\Tr(R_0D_K^2)=\Tr D_K^2$.
Diese neutrale Antwort misst die Gesamtnorm, wählt aber keine
Hierarchie oder relative Majoranaphase.

\textbf{Derselbe Block erzeugt Seesaw, Determinante und Gedächtnis.}
Die bestehende bedingte Neutrinorechnung liefert die schwere Skala
$m_*=M_3$ in eV und die normierte Matrix $\widehat Y$ mit Einheit
$\sqrt{\mathrm{eV}}$. Setze
\begin{equation}
H=D_{M_R/M_3},\qquad
\widehat V=\diag(\widehat Y^{\mathsf T},\widehat Y^\dagger),\qquad
z=E/m_*.
\end{equation}
Der ursprüngliche endliche Neutrinooperator besitzt nach der entsprechenden
Umsortierung die dimensionshomogene Form
\begin{equation}
\mathcal D_{\rm phys}=
\begin{pmatrix}
0&\sqrt{m_*}\,\widehat V\\
\sqrt{m_*}\,\widehat V^\dagger&m_*H
\end{pmatrix}.
\end{equation}
Alle Einträge tragen hier die Einheit eV. Die exakte Elimination des
schweren Blocks liefert
\begin{align}
\Sigma(z)&=\widehat V(zI-H)^{-1}\widehat V^\dagger,\\
\Sigma(0)&=\begin{pmatrix}0&K_\nu\\K_\nu^\dagger&0\end{pmatrix},
&K_\nu&=-\widehat Y^{\mathsf T}(M_R/M_3)^{-1}\widehat Y,\\
\det(EI-\mathcal D_{\rm phys})
&=\det(EI-m_*H)\,
\det\bigl(EI-\Sigma(E/m_*)\bigr).
\end{align}
Die Determinantenidentität gilt außerhalb der schweren Pole und setzt sich
als entsprechende algebraische Identität fort. Sie ist eine Aussage über
diesen endlichen quadratischen Block. Eine vollständige Majorana-
Pfaffianwirkung oder die ursprüngliche elektromagnetische relative
Determinante ist damit noch nicht identifiziert.

Ist $Hu_j=\lambda_j u_j$, so gilt mit positiven semidefiniten Residuen
\begin{equation}
\Sigma(z)=\sum_j\frac{\mathcal R_j}{z-\lambda_j},\qquad
\mathcal R_j=(\widehat V u_j)(\widehat V u_j)^\dagger\ge0,\qquad
\mathcal M(u)=\sum_j e^{-iu\lambda_j}\mathcal R_j.
\end{equation}
Die formale retardierte Fourierantwort bezüglich $z$ enthält
$-i\Theta(u)\mathcal M(u)$. Bei einer physischen Zeitdeutung wäre
$u=m_*t$ und die Zeitkern-Normierung entsprechend $m_*\mathcal M(m_*t)$;
eine solche physische Clockidentifikation wird hier nicht hergeleitet.
Außerdem erzeugen anfängliche schwere Zustände eine eigene inhomogene
Antwort, die beim Ausblenden mitgeführt werden muss.

Für den vorhandenen bedingten Neutrinozweig besitzt die Matrixantwort sechs
verschiedene Pole
\begin{equation}
\lambda_j=\pm1,\quad
\pm1.8848376655\cdot10^{-4},\quad
\pm9.4241883277\cdot10^{-5}.
\end{equation}
Alle sechs Residuen sind von null verschieden und von Rang eins. Somit sind
alle sechs Moden in dieser vollständigen Matrixantwort sichtbar; ihre
exakte endliche lineare Realisierung kann diese Pole nicht weglassen.
Das ist kein allgemeiner unterer Speicherbedarf eines beliebigen
Quanten- oder Universalraummodells. Die direkte Resolventenrechnung und
die getrennte Spektralsumme stimmen relativ besser als $5\cdot10^{-16}$
überein; die statische Quellenantwort stimmt mit der bestehenden leichten
Seesaw-Matrix bis auf $4.1\cdot10^{-17}$ eV überein.
Diese Zahlen prüfen die gemeinsame Implementierung, nicht eine neue
Messvorhersage.

\picturetext{Die schweren Zustände sind wie eine verborgene Mechanik hinter
einem sichtbaren Hebel. Bewegt man den Hebel sehr langsam, spürt man nur
ihren wirksamen Widerstand: Das ist hier die leichte Massenmatrix. Ändert
man die Anregung, antwortet die verborgene Mechanik mit ihren eigenen
Moden weiter. Genau diese Rückwirkung muss eine korrekte Vergröberung
speichern. Masse, Determinantenantwort und Gedächtnis entstehen in dieser
Rechnung aus derselben gekoppelten Struktur.}

\textbf{Was dadurch über die gesuchte gemeinsame Lösung feststeht.}
Der Majoranablock benötigt auf der Ebene dieser inneren Operatorwirkung
keine hinzugefügte fremde Quelle. Die Wahl einer festen geladenen
Einsetzung bleibt jedoch eine echte Wahl:
$\omega D_K\omega^{-1}=-D_K$. Ein physischer Massenterm verlangt daher
die passende Higgs-/Phasenrealisierung und dieselbe ursprüngliche
Wirkung, welche auch die Nahtkopplung auswählt. Die hier konstruierte
lineare Abbildung stellt \emph{jede} symmetrische Massenform dar; ihre
Existenz allein bevorzugt die gewünschte Form nicht.
Ebenso wenig bestimmt die allgemeine Schuridentität den besonderen
$\alpha$-Kanal mit Gewicht $5/4$ und Kopplung $c_3^2$.
Der nächste Herkunftsschritt muss deshalb am ursprünglichen gemeinsamen
Variationsproblem ansetzen: Es muss die bereits eingeschränkten
geladenen Einsetzungen, ihre Normierung und die primitive Nahtantwort
als zusammengehörigen stationären Prozess auswählen.

